% Final §3 -- Architecture and implementation (plan §7.1: 2.5 pages).
% Source of truth: docs/ARCHITECTURE.md (§1-§14) and the code it names; README.md §8 (claim -> code).
\section{Architecture and implementation}\label{sec:architecture}

\begin{figure}[t]
  \centering% Fig. 2 (updated for the Final report): cloud architecture of SeatRelay.
% Source of truth: docs/ARCHITECTURE.md §1.2 (Figure A), §3–§7. Needs % SeatRelay report: shared TikZ setup for the hand-drawn figures.
%
% Load once in the report preamble:   % SeatRelay report: shared TikZ setup for the hand-drawn figures.
%
% Load once in the report preamble:   \input{figures/tikz-setup}
% Then place a figure with
%   \begin{figure}[t]\centering\input{figures/fig-architecture}\caption{...}\label{fig:arch}\end{figure}
%
% The figure files contain only a tikzpicture (no figure environment, no caption), so captions
% and labels stay in the report text. All diagram text is set at a fixed 7 pt sans size, so the
% figures have the same physical size whatever the document's base font size is. Widths are
% checked against a single-column A4 layout with 2 cm margins (\textwidth = 17 cm): every
% figure is at most 0.95\textwidth wide.
%
% Colours: the categorical slots of the validated palette used by the result charts
% (report/tools/make_figures.py), as light tints for node fills. Identity never rests on colour
% alone: every node names its service in text.

\usepackage{tikz}
\usetikzlibrary{arrows.meta,positioning,calc,fit,backgrounds,shapes.geometric,decorations.pathreplacing}

\definecolor{srBlue}{HTML}{2A78D6}
\definecolor{srOrange}{HTML}{EB6834}
\definecolor{srGreen}{HTML}{1BAF7A}
\definecolor{srYellow}{HTML}{EDA100}
\definecolor{srPurple}{HTML}{6250D6}
\definecolor{srInk}{HTML}{1F1F1D}
\definecolor{srMuted}{HTML}{52514E}
\definecolor{srGrid}{HTML}{D8D7D2}

\newcommand{\srfont}{\sffamily\fontsize{7}{8.2}\selectfont}
\newcommand{\srfontb}{\sffamily\bfseries\fontsize{7}{8.2}\selectfont}

\tikzset{
  sr/base/.style={draw=srInk!75, line width=0.5pt, rounded corners=1.6pt, align=center,
    inner xsep=2.5pt, inner ysep=2pt, font=\srfont, text=srInk, minimum height=9.5mm},
  % client and edge (browser, CloudFront, S3)
  sr/client/.style={sr/base, fill=srBlue!10},
  % identity, scheduling, operations (Cognito, EventBridge, CloudWatch, SES)
  sr/service/.style={sr/base, fill=srPurple!8},
  % compute (Lambda functions, API Gateway)
  sr/compute/.style={sr/base, fill=srOrange!11},
  % queues (SQS)
  sr/queue/.style={sr/base, fill=srYellow!16},
  % data (DynamoDB)
  sr/data/.style={sr/base, fill=srGreen!11},
  % numbered step marker, placed on a path: node[sr/step]{3}
  sr/step/.style={circle, fill=srInk, text=white, font=\sffamily\bfseries\fontsize{6}{6}\selectfont,
    inner sep=0pt, minimum size=8.4pt},
  % arrows
  sr/flow/.style={draw=srInk, line width=0.55pt, -{Stealth[length=3.6pt,width=2.8pt]}},
  sr/async/.style={sr/flow, dash pattern=on 2.4pt off 1.4pt},
  sr/reply/.style={sr/flow, densely dotted},
  sr/aux/.style={draw=srMuted, line width=0.45pt, dotted, -{Stealth[length=3.2pt,width=2.6pt]}},
  % arrow and annotation labels
  sr/lbl/.style={font=\srfont, text=srInk, align=center, inner sep=1pt},
  sr/note/.style={font=\srfont, text=srMuted, align=center, inner sep=1pt},
}

% Then place a figure with
%   \begin{figure}[t]\centering% Fig. 2 (updated for the Final report): cloud architecture of SeatRelay.
% Source of truth: docs/ARCHITECTURE.md §1.2 (Figure A), §3–§7. Needs \input{figures/tikz-setup}.
% Size: about 15.8 cm x 6.5 cm (fits 0.95\textwidth of A4 with 2 cm margins; height <= 7 cm).
%
% Step numbers follow one write as in Figure 2 of the Preliminary report: (1) assets, (2) sign-in,
% (3) POST through the JWT authorizer to the api function (which records CMD#{id} = ACCEPTED),
% (4) enqueue in the workshop's FIFO group and answer 202, (5) the worker (single writer) runs
% decide(), (6) one TransactWriteItems guarded by the META version, (7) notifications via Streams,
% notifier and SES, (8a/8b) the per-offer timer re-enters the same FIFO group, (9) the browser
% polls GET /commands/{id}. Changed since the Preliminary report: the timer path 8a/8b replaces the
% expiry sweep (its old step 8), which now runs every 5 min only as a safety net (grey, unnumbered);
% notifications leave after the commit from the stream instead of from the worker. Nodes marked
% with an asterisk are new. Solid arrows: requests; dotted black: the 202 reply; dashed: queued
% messages and timers; grey: the safety net and the sweep's index query.
\begin{tikzpicture}[x=1cm, y=1cm]
  % ---------------------------------------------------------------- top band
  \node[sr/client,  minimum width=2.40cm] (cf)    at (1.20,6.00)  {\textbf{CloudFront + S3}\\private bucket, OAC};
  \node[sr/service, minimum width=1.90cm] (cog)   at (3.95,6.00)  {\textbf{Cognito}\\hosted UI};
  \node[sr/service, minimum width=1.90cm] (eb)    at (6.55,6.00)  {\textbf{EventBridge}\\Scheduler, 5 min};
  \node[sr/compute, minimum width=1.30cm] (sweep) at (8.75,6.00)  {\textbf{sweep}\\safety net};
  \node[sr/compute, minimum width=2.30cm] (fwd)   at (11.05,6.00) {\textbf{expiryForwarder}$^{\ast}$\\Lambda};
  \node[sr/queue,   minimum width=2.90cm] (delay) at (14.25,6.00) {\textbf{SQS expiry-delay}$^{\ast}$\\standard queue, one\\timer per offer, $\le$\,900\,s};
  % ---------------------------------------------------------------- request band
  \node[sr/client,  minimum width=1.65cm] (spa)   at (0.825,3.55)  {\textbf{Browser}\\SPA (React)};
  \node[sr/compute, minimum width=2.10cm] (gw)    at (5.82,3.55)  {\textbf{API Gateway}\\HTTP API\\JWT authorizer};
  \node[sr/compute, minimum width=1.30cm] (api)   at (8.00,3.55)  {\textbf{api}\\Lambda};
  \node[sr/queue,   minimum width=2.70cm] (fifo)  at (10.55,3.55) {\textbf{SQS FIFO}\\commands.fifo\\group = workshopId};
  \node[sr/compute, minimum width=2.90cm] (wk)    at (14.25,3.55) {\textbf{worker} Lambda\\tick-first decide()\\TransactWriteItems};
  % ---------------------------------------------------------------- data band
  \node[sr/data, minimum width=9.50cm, minimum height=1.10cm] (ddb) at (10.95,1.30)
    {\textbf{DynamoDB single table} (on-demand, PITR, TTL, Streams)\\
     CMD\#$\{$id$\}$: command results \enspace$\cdot$\enspace WS\#$\{$id$\}$: META with version lock,\\
     live L\# and history H\# items \enspace$\cdot$\enspace sparse DUE index (GSI1)};
  \node[sr/compute, minimum width=1.60cm] (nt)    at (3.80,1.30)  {\textbf{notifier}$^{\ast}$\\Lambda};
  \node[sr/service, minimum width=1.90cm] (ses)   at (0.95,1.30)  {\textbf{Amazon SES}\\e-mail};
  % ---------------------------------------------------------------- operations strip
  \node[sr/service, minimum width=15.70cm, minimum height=0.42cm, inner ysep=1pt, anchor=south west] (cw) at (0,0)
    {\textbf{CloudWatch}: EMF metrics, dashboard, alarms $\to$ SNS
     \enspace$\cdot$\enspace a DLQ behind every queue \enspace$\cdot$\enspace AWS Budgets (credits excluded)};

  % ---------------------------------------------------------------- (1) static assets, (2) sign-in
  \draw[sr/flow] (spa.north -| cf.south) -- node[sr/step, pos=0.45] {1} node[sr/lbl, left=5pt, pos=0.45] {assets} (cf.south);
  \draw[sr/flow, {Stealth[length=3.6pt,width=2.8pt]}-{Stealth[length=3.6pt,width=2.8pt]}]
    ($(spa.north east)+(-0.25,0)$) -- node[sr/step, pos=0.35] {2}
    node[sr/lbl, pos=0.5, right=6pt, align=left] {sign in,\\ID token} ($(cog.south west)+(0.40,0)$);
  % ---------------------------------------------------------------- (3) POST, 202, (9) poll
  \draw[sr/flow]  (spa.east |- 0,3.92) -- node[sr/step, pos=0.07] {3} node[sr/lbl, above, pos=0.58] {POST \ldots/commands} (gw.west |- 0,3.92);
  \draw[sr/reply] (gw.west |- 0,3.55) -- node[sr/lbl, above, pos=0.5] {202 Accepted} (spa.east |- 0,3.55);
  \draw[sr/flow]  (spa.east |- 0,3.18) -- node[sr/step, pos=0.07] {9} node[sr/lbl, above, pos=0.58] {GET /commands/$\{$id$\}$} (gw.west |- 0,3.18);
  \node[sr/note, anchor=north] at ($(spa.east |- 0,3.02)!0.5!(gw.west |- 0,3.02)$) {poll 300\,$\to$\,2000\,ms until\\APPLIED or REJECTED};
  % ---------------------------------------------------------------- (4) CMD record, (5) enqueue, (6) deliver
  \draw[sr/flow] (gw) -- (api);
  \draw[sr/flow] (api.south) -- node[sr/lbl, left=1pt, pos=0.32, align=right] {CMD\#$\{$id$\}$\\put, get} (api.south |- ddb.north);
  \draw[sr/flow] (api) -- node[sr/step] {4} (fifo);
  \draw[sr/flow] (fifo) -- node[sr/step] {5} (wk);
  % ---------------------------------------------------------------- worker reads and commits
  \coordinate (wkR) at ($(wk.south)+(0.85,0)$);
  \coordinate (wkL) at ($(wk.south)+(-0.85,0)$);
  \coordinate (wkT) at ($(wk.north)+(0.75,0)$);
  \draw[sr/flow] (wkR |- ddb.north) -- node[sr/lbl, left=1pt, pos=0.5, align=right] {Query L\#,\\consistent} (wkR);
  \draw[sr/flow] (wkL) -- node[sr/step, pos=0.32] {6} node[sr/lbl, left=3pt, pos=0.32, align=right] {one transaction,\\META version $= v$} (wkL |- ddb.north);
  % ---------------------------------------------------------------- (7) timer, (10) forward, (11) sweep
  \draw[sr/async] (wkT) -- node[sr/step, pos=0.45] {8a} node[sr/lbl, left=5pt, pos=0.45, align=right] {TICK per offer,\\before the commit} (wkT |- delay.south);
  \draw[sr/async] (delay) -- node[sr/step] {8b} (fwd);
  \coordinate (fwdB) at ($(fwd.south)+(-0.05,0)$);
  \draw[sr/async] (fwdB) -- node[sr/lbl, right=2pt, pos=0.45, align=left] {TICK,\\same group} (fwdB |- fifo.north);
  \draw[sr/flow, draw=srMuted] (eb) -- (sweep);
  \coordinate (swB) at ($(sweep.south)+(0.45,0)$);
  \draw[sr/async, draw=srMuted] (swB) -- node[sr/note, right=1pt, pos=0.45, align=left] {TICK if\\overdue} (swB |- fifo.north);
  \coordinate (gap) at ($(api.east)!0.5!(fifo.west)$);
  \draw[sr/aux] (gap |- sweep.south) -- node[sr/note, pos=0.22, left=1pt, align=right] {Query\\DUE index} (gap |- ddb.north);
  % ---------------------------------------------------------------- (12) notifications
  \draw[sr/flow] (ddb.west) -- node[sr/step] {7} node[sr/lbl, above=4pt] {Streams} (nt.east);
  \draw[sr/flow] (nt) -- node[sr/lbl, above] {send} (ses);
\end{tikzpicture}
\caption{...}\label{fig:arch}\end{figure}
%
% The figure files contain only a tikzpicture (no figure environment, no caption), so captions
% and labels stay in the report text. All diagram text is set at a fixed 7 pt sans size, so the
% figures have the same physical size whatever the document's base font size is. Widths are
% checked against a single-column A4 layout with 2 cm margins (\textwidth = 17 cm): every
% figure is at most 0.95\textwidth wide.
%
% Colours: the categorical slots of the validated palette used by the result charts
% (report/tools/make_figures.py), as light tints for node fills. Identity never rests on colour
% alone: every node names its service in text.

\usepackage{tikz}
\usetikzlibrary{arrows.meta,positioning,calc,fit,backgrounds,shapes.geometric,decorations.pathreplacing}

\definecolor{srBlue}{HTML}{2A78D6}
\definecolor{srOrange}{HTML}{EB6834}
\definecolor{srGreen}{HTML}{1BAF7A}
\definecolor{srYellow}{HTML}{EDA100}
\definecolor{srPurple}{HTML}{6250D6}
\definecolor{srInk}{HTML}{1F1F1D}
\definecolor{srMuted}{HTML}{52514E}
\definecolor{srGrid}{HTML}{D8D7D2}

\newcommand{\srfont}{\sffamily\fontsize{7}{8.2}\selectfont}
\newcommand{\srfontb}{\sffamily\bfseries\fontsize{7}{8.2}\selectfont}

\tikzset{
  sr/base/.style={draw=srInk!75, line width=0.5pt, rounded corners=1.6pt, align=center,
    inner xsep=2.5pt, inner ysep=2pt, font=\srfont, text=srInk, minimum height=9.5mm},
  % client and edge (browser, CloudFront, S3)
  sr/client/.style={sr/base, fill=srBlue!10},
  % identity, scheduling, operations (Cognito, EventBridge, CloudWatch, SES)
  sr/service/.style={sr/base, fill=srPurple!8},
  % compute (Lambda functions, API Gateway)
  sr/compute/.style={sr/base, fill=srOrange!11},
  % queues (SQS)
  sr/queue/.style={sr/base, fill=srYellow!16},
  % data (DynamoDB)
  sr/data/.style={sr/base, fill=srGreen!11},
  % numbered step marker, placed on a path: node[sr/step]{3}
  sr/step/.style={circle, fill=srInk, text=white, font=\sffamily\bfseries\fontsize{6}{6}\selectfont,
    inner sep=0pt, minimum size=8.4pt},
  % arrows
  sr/flow/.style={draw=srInk, line width=0.55pt, -{Stealth[length=3.6pt,width=2.8pt]}},
  sr/async/.style={sr/flow, dash pattern=on 2.4pt off 1.4pt},
  sr/reply/.style={sr/flow, densely dotted},
  sr/aux/.style={draw=srMuted, line width=0.45pt, dotted, -{Stealth[length=3.2pt,width=2.6pt]}},
  % arrow and annotation labels
  sr/lbl/.style={font=\srfont, text=srInk, align=center, inner sep=1pt},
  sr/note/.style={font=\srfont, text=srMuted, align=center, inner sep=1pt},
}
.
% Size: about 15.8 cm x 6.5 cm (fits 0.95\textwidth of A4 with 2 cm margins; height <= 7 cm).
%
% Step numbers follow one write as in Figure 2 of the Preliminary report: (1) assets, (2) sign-in,
% (3) POST through the JWT authorizer to the api function (which records CMD#{id} = ACCEPTED),
% (4) enqueue in the workshop's FIFO group and answer 202, (5) the worker (single writer) runs
% decide(), (6) one TransactWriteItems guarded by the META version, (7) notifications via Streams,
% notifier and SES, (8a/8b) the per-offer timer re-enters the same FIFO group, (9) the browser
% polls GET /commands/{id}. Changed since the Preliminary report: the timer path 8a/8b replaces the
% expiry sweep (its old step 8), which now runs every 5 min only as a safety net (grey, unnumbered);
% notifications leave after the commit from the stream instead of from the worker. Nodes marked
% with an asterisk are new. Solid arrows: requests; dotted black: the 202 reply; dashed: queued
% messages and timers; grey: the safety net and the sweep's index query.
\begin{tikzpicture}[x=1cm, y=1cm]
  % ---------------------------------------------------------------- top band
  \node[sr/client,  minimum width=2.40cm] (cf)    at (1.20,6.00)  {\textbf{CloudFront + S3}\\private bucket, OAC};
  \node[sr/service, minimum width=1.90cm] (cog)   at (3.95,6.00)  {\textbf{Cognito}\\hosted UI};
  \node[sr/service, minimum width=1.90cm] (eb)    at (6.55,6.00)  {\textbf{EventBridge}\\Scheduler, 5 min};
  \node[sr/compute, minimum width=1.30cm] (sweep) at (8.75,6.00)  {\textbf{sweep}\\safety net};
  \node[sr/compute, minimum width=2.30cm] (fwd)   at (11.05,6.00) {\textbf{expiryForwarder}$^{\ast}$\\Lambda};
  \node[sr/queue,   minimum width=2.90cm] (delay) at (14.25,6.00) {\textbf{SQS expiry-delay}$^{\ast}$\\standard queue, one\\timer per offer, $\le$\,900\,s};
  % ---------------------------------------------------------------- request band
  \node[sr/client,  minimum width=1.65cm] (spa)   at (0.825,3.55)  {\textbf{Browser}\\SPA (React)};
  \node[sr/compute, minimum width=2.10cm] (gw)    at (5.82,3.55)  {\textbf{API Gateway}\\HTTP API\\JWT authorizer};
  \node[sr/compute, minimum width=1.30cm] (api)   at (8.00,3.55)  {\textbf{api}\\Lambda};
  \node[sr/queue,   minimum width=2.70cm] (fifo)  at (10.55,3.55) {\textbf{SQS FIFO}\\commands.fifo\\group = workshopId};
  \node[sr/compute, minimum width=2.90cm] (wk)    at (14.25,3.55) {\textbf{worker} Lambda\\tick-first decide()\\TransactWriteItems};
  % ---------------------------------------------------------------- data band
  \node[sr/data, minimum width=9.50cm, minimum height=1.10cm] (ddb) at (10.95,1.30)
    {\textbf{DynamoDB single table} (on-demand, PITR, TTL, Streams)\\
     CMD\#$\{$id$\}$: command results \enspace$\cdot$\enspace WS\#$\{$id$\}$: META with version lock,\\
     live L\# and history H\# items \enspace$\cdot$\enspace sparse DUE index (GSI1)};
  \node[sr/compute, minimum width=1.60cm] (nt)    at (3.80,1.30)  {\textbf{notifier}$^{\ast}$\\Lambda};
  \node[sr/service, minimum width=1.90cm] (ses)   at (0.95,1.30)  {\textbf{Amazon SES}\\e-mail};
  % ---------------------------------------------------------------- operations strip
  \node[sr/service, minimum width=15.70cm, minimum height=0.42cm, inner ysep=1pt, anchor=south west] (cw) at (0,0)
    {\textbf{CloudWatch}: EMF metrics, dashboard, alarms $\to$ SNS
     \enspace$\cdot$\enspace a DLQ behind every queue \enspace$\cdot$\enspace AWS Budgets (credits excluded)};

  % ---------------------------------------------------------------- (1) static assets, (2) sign-in
  \draw[sr/flow] (spa.north -| cf.south) -- node[sr/step, pos=0.45] {1} node[sr/lbl, left=5pt, pos=0.45] {assets} (cf.south);
  \draw[sr/flow, {Stealth[length=3.6pt,width=2.8pt]}-{Stealth[length=3.6pt,width=2.8pt]}]
    ($(spa.north east)+(-0.25,0)$) -- node[sr/step, pos=0.35] {2}
    node[sr/lbl, pos=0.5, right=6pt, align=left] {sign in,\\ID token} ($(cog.south west)+(0.40,0)$);
  % ---------------------------------------------------------------- (3) POST, 202, (9) poll
  \draw[sr/flow]  (spa.east |- 0,3.92) -- node[sr/step, pos=0.07] {3} node[sr/lbl, above, pos=0.58] {POST \ldots/commands} (gw.west |- 0,3.92);
  \draw[sr/reply] (gw.west |- 0,3.55) -- node[sr/lbl, above, pos=0.5] {202 Accepted} (spa.east |- 0,3.55);
  \draw[sr/flow]  (spa.east |- 0,3.18) -- node[sr/step, pos=0.07] {9} node[sr/lbl, above, pos=0.58] {GET /commands/$\{$id$\}$} (gw.west |- 0,3.18);
  \node[sr/note, anchor=north] at ($(spa.east |- 0,3.02)!0.5!(gw.west |- 0,3.02)$) {poll 300\,$\to$\,2000\,ms until\\APPLIED or REJECTED};
  % ---------------------------------------------------------------- (4) CMD record, (5) enqueue, (6) deliver
  \draw[sr/flow] (gw) -- (api);
  \draw[sr/flow] (api.south) -- node[sr/lbl, left=1pt, pos=0.32, align=right] {CMD\#$\{$id$\}$\\put, get} (api.south |- ddb.north);
  \draw[sr/flow] (api) -- node[sr/step] {4} (fifo);
  \draw[sr/flow] (fifo) -- node[sr/step] {5} (wk);
  % ---------------------------------------------------------------- worker reads and commits
  \coordinate (wkR) at ($(wk.south)+(0.85,0)$);
  \coordinate (wkL) at ($(wk.south)+(-0.85,0)$);
  \coordinate (wkT) at ($(wk.north)+(0.75,0)$);
  \draw[sr/flow] (wkR |- ddb.north) -- node[sr/lbl, left=1pt, pos=0.5, align=right] {Query L\#,\\consistent} (wkR);
  \draw[sr/flow] (wkL) -- node[sr/step, pos=0.32] {6} node[sr/lbl, left=3pt, pos=0.32, align=right] {one transaction,\\META version $= v$} (wkL |- ddb.north);
  % ---------------------------------------------------------------- (7) timer, (10) forward, (11) sweep
  \draw[sr/async] (wkT) -- node[sr/step, pos=0.45] {8a} node[sr/lbl, left=5pt, pos=0.45, align=right] {TICK per offer,\\before the commit} (wkT |- delay.south);
  \draw[sr/async] (delay) -- node[sr/step] {8b} (fwd);
  \coordinate (fwdB) at ($(fwd.south)+(-0.05,0)$);
  \draw[sr/async] (fwdB) -- node[sr/lbl, right=2pt, pos=0.45, align=left] {TICK,\\same group} (fwdB |- fifo.north);
  \draw[sr/flow, draw=srMuted] (eb) -- (sweep);
  \coordinate (swB) at ($(sweep.south)+(0.45,0)$);
  \draw[sr/async, draw=srMuted] (swB) -- node[sr/note, right=1pt, pos=0.45, align=left] {TICK if\\overdue} (swB |- fifo.north);
  \coordinate (gap) at ($(api.east)!0.5!(fifo.west)$);
  \draw[sr/aux] (gap |- sweep.south) -- node[sr/note, pos=0.22, left=1pt, align=right] {Query\\DUE index} (gap |- ddb.north);
  % ---------------------------------------------------------------- (12) notifications
  \draw[sr/flow] (ddb.west) -- node[sr/step] {7} node[sr/lbl, above=4pt] {Streams} (nt.east);
  \draw[sr/flow] (nt) -- node[sr/lbl, above] {send} (ses);
\end{tikzpicture}

  \caption{Cloud architecture and request path. (1)~SPA from a private S3 bucket via CloudFront.
  (2)~Sign-in on the Cognito hosted page. (3)~\code{POST /workshops/\{id\}/commands}
  passes the JWT authorizer; (4)~\code{api} writes \code{CMD\#\{commandId\}} = ACCEPTED and
  (5)~enqueues the command in the workshop's FIFO group, answering 202. (6)~The worker, the only
  writer of a workshop, runs \code{decide()}, (7)~sends one TICK per new offer to the delay queue and
  (8)~commits one transaction guarded by the workshop version. (9)~The browser polls
  \code{GET /commands/\{id\}}. (10)~Due timers are forwarded into the FIFO group; (11)~a 5-minute
  sweep re-ticks overdue workshops (safety net); (12)~notifications go out via Streams and SES.
  Solid: requests; dotted: the 202 reply; dashed: queued messages; grey dotted: index query.}
  \label{fig:arch}
\end{figure}

\subsection{Components and layering}
SeatRelay runs in one AWS region (\code{ap-southeast-1}) as three CloudFormation stacks: identity
(Cognito user pool, Simple Email Service (SES) identity, budgets), application (one per stage: \code{dev} for experiments,
\code{pilot} for real users) and web (private S3 bucket behind CloudFront with origin access
control~\cite{aws-cloudfront-oac}). Fig.~\ref{fig:arch} shows the application stack. Users sign in on the Cognito hosted page (authorization code flow with Proof Key for Code Exchange,
PKCE), and the HTTP API validates the JSON Web Token's (JWT) signature, issuer and audience before
any function runs~\cite{aws-apigw-jwt}. Five Lambda functions (arm64, Node.js~22) do the work:
\code{api}, \code{worker}, \code{expiryForwarder}, \code{sweep} and \code{notifier}. Two Simple Queue
Service (SQS) queues carry commands and timers: a first-in-first-out (FIFO) command queue and a
standard delay queue. State lives in one DynamoDB table (on-demand capacity, point-in-time recovery
(PITR), Streams, time-to-live (TTL)). We chose pay-per-request managed services because registration
traffic is bursty (opening of registration, deadlines) and idle most of the time, so idle cost is
close to zero (\S\ref{sec:business}), and because they run across AZs without servers to patch.

All business rules are in a pure engine, \code{decide(snapshot, command, now)} in
\code{packages/core}, which returns the writes, notifications, timers and backfill records of one
decision and performs no I/O; thin adapters map its changes to DynamoDB items or PostgreSQL rows. The
same engine runs in the cloud, in the on-prem reference and in the replay simulator
(\S\ref{sec:mechanism}).

\subsection{Write path, data model and the single-writer lock}\label{sec:writepath}
Reads, organisation and workshop set-up and attendance are synchronous. Every change to a
workshop's live state is a \emph{command} with a client-generated \code{commandId}. The API writes
the command record \code{CMD\#\{commandId\}} with status ACCEPTED (conditional on absence; a repeat
returns the existing record), enqueues the command on the FIFO queue with message group = workshop
and deduplication id = \code{commandId}~\cite{aws-sqs-exactly-once}, and answers 202. If the
enqueue fails, the record is deleted and the API answers 503, so no browser polls an orphan. SQS
FIFO hands the messages of one group to the function one at a time and in
order~\cite{aws-sqs-fifo-terms,aws-lambda-sqs-scaling}; each workshop therefore has exactly one
serial writer, while different workshops are processed in parallel (worker concurrency 5).

The table keeps a workshop in one partition \code{WS\#\{id\}}: the \emph{live} items that a decision
needs under prefix \code{L\#} (workshop META item, sessions, live registrations and offers, pending
backfills, a due-date marker) and settled \emph{history} under \code{H\#} (archived registrations
and offers, backfill records, attendance). The worker reads only \code{L\#} with a strongly
consistent query, so its snapshot is bounded by capacity and queue length, not by the workshop's
age. It runs \code{decide()}: expire what is due (\S\ref{sec:timers}), apply the command or reject
it with a deterministic reason, allocate free seats in queue order, and recompute the next due
time. It then commits \emph{one} \code{TransactWriteItems}~\cite{aws-ddb-transactwriteitems} in
which the META item is conditioned on \code{version = }the snapshot's version and incremented. This
optimistic lock covers the whole workshop and is a second line of defence behind the single writer;
a version conflict reloads and decides again (at most three times). The command record is
conditioned on status ACCEPTED in the same transaction, so a command's effects and its settlement
commit together or not at all, and a redelivered or duplicated command fails the condition and is
acknowledged without effect (C3). A transaction holds at most 100 actions, so a decision is capped
at 12 offers, 12 expiries and 90 items; remaining work continues as an internal REALLOCATE command.
A deterministic rejection is a successful settlement (REJECTED) and never throws; an infrastructure
error reports the failed message and all later ones of the batch for redelivery, as FIFO sources
require~\cite{aws-lambda-sqs-errors}. They return after the 120\,s visibility timeout; three failed
receives move a message to a dead-letter queue (DLQ) that raises an alarm. Command records expire
after 30 days and notifications after 90 days (TTL).

\subsection{Rules R1--R7 and constraints C1--C7}
Table~\ref{tab:rules} maps each rule and constraint of the Preliminary report to its enforcement.
Every rule and every rejection reason has a unit test, and a property-based
test~\cite{claessen2000quickcheck,fastcheck} applies random command sequences with random clock
jumps and checks the invariants after every step (\S\ref{sec:results}). Organisers can run the same
invariant check on a live workshop from the dashboard (\code{verify}).

\begin{table}[t]
\centering\footnotesize
\caption{Rules and constraints and where they are enforced (functions in \code{packages/core} unless
stated; full map in \code{docs/ARCHITECTURE.md} \S9--10).}\label{tab:rules}
\begin{tabularx}{\textwidth}{@{}l P{6.6cm} L@{}}
\toprule
 & Statement & Enforced by\\
\midrule
R1 & One live registration per account per workshop. & \code{register()} rejects ALREADY\_REGISTERED under the single writer and version lock; mirrored by a marker item per user.\\
R2 & Queue order = FIFO arrival order = \code{seq} assigned by the worker. & \code{seq = META.seq + 1}, written only by the workshop's single writer.\\
R3 & Allocate in \code{seq} order; offer the first acceptable session, in preference order, that has a free seat. & \code{allocate()}, \code{firstAllocatableSession()}.\\
R4 & $\text{expiresAt} = \min(\text{now} + \text{holdSeconds}, \text{cutoffAt})$, cutoffAt = start $-$ 30\,min; no offer after cut-off. & \code{offerExpiry()}, \code{tickFirst()}; per-offer timers and the sweep (\S\ref{sec:timers}).\\
R5 & Decline or expiry drops that session; back to QUEUED with the same \code{seq}, or ENDED if none is left. & \code{declineOffer()}, \code{expireOffer()}; \code{allocate()} in the same decision.\\
R6 & Cancellation exits completely; the seat is relayed at once; re-registering gets a new \code{seq}. & \code{cancelRegistration()}; relay in the same transaction.\\
R7 & Capacity below confirmed + held is rejected; an increase allocates at once; cancelling a session revokes its offers and re-queues its confirmed participants. & \code{setCapacity()}, \code{cancelSession()} and \code{drainCancellations()}.\\
\midrule
C1 & No over-capacity: confirmed + held $\leq$ capacity. & Offers only when free $> 0$; counters change only in \code{decide()} and commit atomically; \code{checkInvariants()}.\\
C2 & At most one OFFERED or CONFIRMED seat per account per workshop. & R1 plus one current offer per registration; only QUEUED registrations receive offers.\\
C3 & A command takes effect at most once. & Client \code{commandId}; conditional record; FIFO deduplication; record condition inside the effect transaction.\\
C4 & \pending{C4 wording from the Preliminary report} & Candidates: per-workshop ordering (FIFO group, R2); cut-off (R4); notifications written in the state change's transaction; audit trail (command \code{acceptedAt}/\code{settledAt}, backfill records with latency, history items).\\
C5 & Expiry is decided by server time when a command is processed; a late tick cannot invalidate a newer offer (reworded, \S\ref{sec:timers}). & \code{tickFirst()} before every command; ACCEPT checks \code{expiresAt}; ticks never expire an offer by id.\\
C6 & \pending{C6 wording from the Preliminary report} & Same candidates as C4.\\
C7 & Organiser actions only by members of the workshop's organisation. & Membership checks in the API (403/404) and again in the worker; command results visible only to issuer and organisers; OFFER\_NOT\_YOURS.\\
\bottomrule
\end{tabularx}
\end{table}

\subsection{Responses to the TA's feedback}\label{sec:timers}
\paragraph{(1) Market size.} Added in \S\ref{sec:motivation} (Table~\ref{tab:market}).

\paragraph{(2) Per-offer timers instead of sweep-bounded expiry.}
In the Preliminary design a periodic sweep expired offers, so the sweep interval bounded the
backfill latency. Each offer now has its own timer (Fig.~\ref{fig:seq}b). SQS FIFO queues do not
support timers on individual messages, and a message timer is at most 15
minutes~\cite{aws-sqs-timers}. The worker therefore sends one TICK per new offer to a
\emph{standard} delay queue with $\text{DelaySeconds} = \min(900, \lceil \text{expiresAt} -
\text{now} \rceil)$, \emph{before} its commit (a timer whose commit fails finds nothing to expire).
When the TICK becomes visible, the \code{expiryForwarder} moves it into the workshop's FIFO group
(deduplication id \code{tick\#offerId\#expiresAt}), so expiries take the same ordered single-writer
path as user commands; a TICK that is still early, as for holds longer than 15 minutes, goes back to
the delay queue with the remaining delay. A TICK is only a wake-up call: every decision begins by
expiring all offers with $\text{expiresAt} \leq \text{now}$ (\emph{tick-first}). Hence late or
duplicate ticks are no-ops, an ACCEPT processed after \code{expiresAt} is rejected whether or not
its timer has arrived, and a late tick cannot touch a newer offer; this is the reworded C5. The
5-minute sweep stays as a safety net: each commit maintains the workshop's next due time on a sparse
index, and the sweep enqueues a TICK for every workshop whose due time has passed (lost timers,
session cut-offs, lost continuations). A metric counts expiries that only the sweep performed
(\code{SweepLateExpiries}, expected 0). The sweep interval therefore no longer bounds the backfill
latency, which is the delay-queue visibility (1\,s granularity) plus the forwarder, one FIFO
enqueue and one decision; a relay caused by a decline or cancellation happens inside the releasing
transaction. We considered one-time EventBridge Scheduler schedules per
offer~\cite{aws-eventbridge-scheduler-types}, which have no 15-minute limit, but they add a service,
a schedule to create and delete per offer and an execution role, while default holds (900\,s) fit
one SQS timer and longer holds take extra hops.

\paragraph{(3) How the browser learns that a 202 has settled: polling.}
The 202 response carries \code{href} and \code{pollAfterMs} = 300 (Fig.~\ref{fig:seq}a). The browser
polls \code{GET /commands/\{id\}}, a strongly consistent read, after 300, 500 and 1000\,ms and then
every 2000\,ms; after 30\,s it polls every 5\,s and shows a \enquote{still processing} note. The
terminal state is APPLIED with a result or REJECTED with a reason; the result contains the
participant's registration summary (status, queue position, offer expiry), so the interface updates
from the poll alone and separates \enquote{request received} from \enquote{seat confirmed}. A retry
reuses the \code{commandId}. Changes caused by other people, such as an offer created by a backfill,
reach the status page through a 10\,s refresh while a registration is QUEUED or OFFERED and the tab
is visible. We did not use WebSockets: at pilot scale they would add an API type, a connection
table and a failure path for little gain, and polling is bounded (about
\CostAssumeReadsPerCommand{} reads per command in the cost model). We report the server settle
latency (\code{settledAt} $-$ \code{acceptedAt}, both on the command record) separately from the
client-observed latency, which includes up to one polling interval.

\begin{figure}[t]
  \centering% Fig. 3 (plan §7.1): request sequences, two compact diagrams side by side.
% (a) command path: 202 -> poll -> settled; (b) expiry relay: offer -> timer -> tick -> backfill.
% Source of truth: docs/ARCHITECTURE.md §3.2 (Figure B) and §5.3 (Figure C); latency
% definitions from ARCHITECTURE.md §3 ("Latency definitions used everywhere").
% Needs % SeatRelay report: shared TikZ setup for the hand-drawn figures.
%
% Load once in the report preamble:   % SeatRelay report: shared TikZ setup for the hand-drawn figures.
%
% Load once in the report preamble:   \input{figures/tikz-setup}
% Then place a figure with
%   \begin{figure}[t]\centering\input{figures/fig-architecture}\caption{...}\label{fig:arch}\end{figure}
%
% The figure files contain only a tikzpicture (no figure environment, no caption), so captions
% and labels stay in the report text. All diagram text is set at a fixed 7 pt sans size, so the
% figures have the same physical size whatever the document's base font size is. Widths are
% checked against a single-column A4 layout with 2 cm margins (\textwidth = 17 cm): every
% figure is at most 0.95\textwidth wide.
%
% Colours: the categorical slots of the validated palette used by the result charts
% (report/tools/make_figures.py), as light tints for node fills. Identity never rests on colour
% alone: every node names its service in text.

\usepackage{tikz}
\usetikzlibrary{arrows.meta,positioning,calc,fit,backgrounds,shapes.geometric,decorations.pathreplacing}

\definecolor{srBlue}{HTML}{2A78D6}
\definecolor{srOrange}{HTML}{EB6834}
\definecolor{srGreen}{HTML}{1BAF7A}
\definecolor{srYellow}{HTML}{EDA100}
\definecolor{srPurple}{HTML}{6250D6}
\definecolor{srInk}{HTML}{1F1F1D}
\definecolor{srMuted}{HTML}{52514E}
\definecolor{srGrid}{HTML}{D8D7D2}

\newcommand{\srfont}{\sffamily\fontsize{7}{8.2}\selectfont}
\newcommand{\srfontb}{\sffamily\bfseries\fontsize{7}{8.2}\selectfont}

\tikzset{
  sr/base/.style={draw=srInk!75, line width=0.5pt, rounded corners=1.6pt, align=center,
    inner xsep=2.5pt, inner ysep=2pt, font=\srfont, text=srInk, minimum height=9.5mm},
  % client and edge (browser, CloudFront, S3)
  sr/client/.style={sr/base, fill=srBlue!10},
  % identity, scheduling, operations (Cognito, EventBridge, CloudWatch, SES)
  sr/service/.style={sr/base, fill=srPurple!8},
  % compute (Lambda functions, API Gateway)
  sr/compute/.style={sr/base, fill=srOrange!11},
  % queues (SQS)
  sr/queue/.style={sr/base, fill=srYellow!16},
  % data (DynamoDB)
  sr/data/.style={sr/base, fill=srGreen!11},
  % numbered step marker, placed on a path: node[sr/step]{3}
  sr/step/.style={circle, fill=srInk, text=white, font=\sffamily\bfseries\fontsize{6}{6}\selectfont,
    inner sep=0pt, minimum size=8.4pt},
  % arrows
  sr/flow/.style={draw=srInk, line width=0.55pt, -{Stealth[length=3.6pt,width=2.8pt]}},
  sr/async/.style={sr/flow, dash pattern=on 2.4pt off 1.4pt},
  sr/reply/.style={sr/flow, densely dotted},
  sr/aux/.style={draw=srMuted, line width=0.45pt, dotted, -{Stealth[length=3.2pt,width=2.6pt]}},
  % arrow and annotation labels
  sr/lbl/.style={font=\srfont, text=srInk, align=center, inner sep=1pt},
  sr/note/.style={font=\srfont, text=srMuted, align=center, inner sep=1pt},
}

% Then place a figure with
%   \begin{figure}[t]\centering% Fig. 2 (updated for the Final report): cloud architecture of SeatRelay.
% Source of truth: docs/ARCHITECTURE.md §1.2 (Figure A), §3–§7. Needs \input{figures/tikz-setup}.
% Size: about 15.8 cm x 6.5 cm (fits 0.95\textwidth of A4 with 2 cm margins; height <= 7 cm).
%
% Step numbers follow one write as in Figure 2 of the Preliminary report: (1) assets, (2) sign-in,
% (3) POST through the JWT authorizer to the api function (which records CMD#{id} = ACCEPTED),
% (4) enqueue in the workshop's FIFO group and answer 202, (5) the worker (single writer) runs
% decide(), (6) one TransactWriteItems guarded by the META version, (7) notifications via Streams,
% notifier and SES, (8a/8b) the per-offer timer re-enters the same FIFO group, (9) the browser
% polls GET /commands/{id}. Changed since the Preliminary report: the timer path 8a/8b replaces the
% expiry sweep (its old step 8), which now runs every 5 min only as a safety net (grey, unnumbered);
% notifications leave after the commit from the stream instead of from the worker. Nodes marked
% with an asterisk are new. Solid arrows: requests; dotted black: the 202 reply; dashed: queued
% messages and timers; grey: the safety net and the sweep's index query.
\begin{tikzpicture}[x=1cm, y=1cm]
  % ---------------------------------------------------------------- top band
  \node[sr/client,  minimum width=2.40cm] (cf)    at (1.20,6.00)  {\textbf{CloudFront + S3}\\private bucket, OAC};
  \node[sr/service, minimum width=1.90cm] (cog)   at (3.95,6.00)  {\textbf{Cognito}\\hosted UI};
  \node[sr/service, minimum width=1.90cm] (eb)    at (6.55,6.00)  {\textbf{EventBridge}\\Scheduler, 5 min};
  \node[sr/compute, minimum width=1.30cm] (sweep) at (8.75,6.00)  {\textbf{sweep}\\safety net};
  \node[sr/compute, minimum width=2.30cm] (fwd)   at (11.05,6.00) {\textbf{expiryForwarder}$^{\ast}$\\Lambda};
  \node[sr/queue,   minimum width=2.90cm] (delay) at (14.25,6.00) {\textbf{SQS expiry-delay}$^{\ast}$\\standard queue, one\\timer per offer, $\le$\,900\,s};
  % ---------------------------------------------------------------- request band
  \node[sr/client,  minimum width=1.65cm] (spa)   at (0.825,3.55)  {\textbf{Browser}\\SPA (React)};
  \node[sr/compute, minimum width=2.10cm] (gw)    at (5.82,3.55)  {\textbf{API Gateway}\\HTTP API\\JWT authorizer};
  \node[sr/compute, minimum width=1.30cm] (api)   at (8.00,3.55)  {\textbf{api}\\Lambda};
  \node[sr/queue,   minimum width=2.70cm] (fifo)  at (10.55,3.55) {\textbf{SQS FIFO}\\commands.fifo\\group = workshopId};
  \node[sr/compute, minimum width=2.90cm] (wk)    at (14.25,3.55) {\textbf{worker} Lambda\\tick-first decide()\\TransactWriteItems};
  % ---------------------------------------------------------------- data band
  \node[sr/data, minimum width=9.50cm, minimum height=1.10cm] (ddb) at (10.95,1.30)
    {\textbf{DynamoDB single table} (on-demand, PITR, TTL, Streams)\\
     CMD\#$\{$id$\}$: command results \enspace$\cdot$\enspace WS\#$\{$id$\}$: META with version lock,\\
     live L\# and history H\# items \enspace$\cdot$\enspace sparse DUE index (GSI1)};
  \node[sr/compute, minimum width=1.60cm] (nt)    at (3.80,1.30)  {\textbf{notifier}$^{\ast}$\\Lambda};
  \node[sr/service, minimum width=1.90cm] (ses)   at (0.95,1.30)  {\textbf{Amazon SES}\\e-mail};
  % ---------------------------------------------------------------- operations strip
  \node[sr/service, minimum width=15.70cm, minimum height=0.42cm, inner ysep=1pt, anchor=south west] (cw) at (0,0)
    {\textbf{CloudWatch}: EMF metrics, dashboard, alarms $\to$ SNS
     \enspace$\cdot$\enspace a DLQ behind every queue \enspace$\cdot$\enspace AWS Budgets (credits excluded)};

  % ---------------------------------------------------------------- (1) static assets, (2) sign-in
  \draw[sr/flow] (spa.north -| cf.south) -- node[sr/step, pos=0.45] {1} node[sr/lbl, left=5pt, pos=0.45] {assets} (cf.south);
  \draw[sr/flow, {Stealth[length=3.6pt,width=2.8pt]}-{Stealth[length=3.6pt,width=2.8pt]}]
    ($(spa.north east)+(-0.25,0)$) -- node[sr/step, pos=0.35] {2}
    node[sr/lbl, pos=0.5, right=6pt, align=left] {sign in,\\ID token} ($(cog.south west)+(0.40,0)$);
  % ---------------------------------------------------------------- (3) POST, 202, (9) poll
  \draw[sr/flow]  (spa.east |- 0,3.92) -- node[sr/step, pos=0.07] {3} node[sr/lbl, above, pos=0.58] {POST \ldots/commands} (gw.west |- 0,3.92);
  \draw[sr/reply] (gw.west |- 0,3.55) -- node[sr/lbl, above, pos=0.5] {202 Accepted} (spa.east |- 0,3.55);
  \draw[sr/flow]  (spa.east |- 0,3.18) -- node[sr/step, pos=0.07] {9} node[sr/lbl, above, pos=0.58] {GET /commands/$\{$id$\}$} (gw.west |- 0,3.18);
  \node[sr/note, anchor=north] at ($(spa.east |- 0,3.02)!0.5!(gw.west |- 0,3.02)$) {poll 300\,$\to$\,2000\,ms until\\APPLIED or REJECTED};
  % ---------------------------------------------------------------- (4) CMD record, (5) enqueue, (6) deliver
  \draw[sr/flow] (gw) -- (api);
  \draw[sr/flow] (api.south) -- node[sr/lbl, left=1pt, pos=0.32, align=right] {CMD\#$\{$id$\}$\\put, get} (api.south |- ddb.north);
  \draw[sr/flow] (api) -- node[sr/step] {4} (fifo);
  \draw[sr/flow] (fifo) -- node[sr/step] {5} (wk);
  % ---------------------------------------------------------------- worker reads and commits
  \coordinate (wkR) at ($(wk.south)+(0.85,0)$);
  \coordinate (wkL) at ($(wk.south)+(-0.85,0)$);
  \coordinate (wkT) at ($(wk.north)+(0.75,0)$);
  \draw[sr/flow] (wkR |- ddb.north) -- node[sr/lbl, left=1pt, pos=0.5, align=right] {Query L\#,\\consistent} (wkR);
  \draw[sr/flow] (wkL) -- node[sr/step, pos=0.32] {6} node[sr/lbl, left=3pt, pos=0.32, align=right] {one transaction,\\META version $= v$} (wkL |- ddb.north);
  % ---------------------------------------------------------------- (7) timer, (10) forward, (11) sweep
  \draw[sr/async] (wkT) -- node[sr/step, pos=0.45] {8a} node[sr/lbl, left=5pt, pos=0.45, align=right] {TICK per offer,\\before the commit} (wkT |- delay.south);
  \draw[sr/async] (delay) -- node[sr/step] {8b} (fwd);
  \coordinate (fwdB) at ($(fwd.south)+(-0.05,0)$);
  \draw[sr/async] (fwdB) -- node[sr/lbl, right=2pt, pos=0.45, align=left] {TICK,\\same group} (fwdB |- fifo.north);
  \draw[sr/flow, draw=srMuted] (eb) -- (sweep);
  \coordinate (swB) at ($(sweep.south)+(0.45,0)$);
  \draw[sr/async, draw=srMuted] (swB) -- node[sr/note, right=1pt, pos=0.45, align=left] {TICK if\\overdue} (swB |- fifo.north);
  \coordinate (gap) at ($(api.east)!0.5!(fifo.west)$);
  \draw[sr/aux] (gap |- sweep.south) -- node[sr/note, pos=0.22, left=1pt, align=right] {Query\\DUE index} (gap |- ddb.north);
  % ---------------------------------------------------------------- (12) notifications
  \draw[sr/flow] (ddb.west) -- node[sr/step] {7} node[sr/lbl, above=4pt] {Streams} (nt.east);
  \draw[sr/flow] (nt) -- node[sr/lbl, above] {send} (ses);
\end{tikzpicture}
\caption{...}\label{fig:arch}\end{figure}
%
% The figure files contain only a tikzpicture (no figure environment, no caption), so captions
% and labels stay in the report text. All diagram text is set at a fixed 7 pt sans size, so the
% figures have the same physical size whatever the document's base font size is. Widths are
% checked against a single-column A4 layout with 2 cm margins (\textwidth = 17 cm): every
% figure is at most 0.95\textwidth wide.
%
% Colours: the categorical slots of the validated palette used by the result charts
% (report/tools/make_figures.py), as light tints for node fills. Identity never rests on colour
% alone: every node names its service in text.

\usepackage{tikz}
\usetikzlibrary{arrows.meta,positioning,calc,fit,backgrounds,shapes.geometric,decorations.pathreplacing}

\definecolor{srBlue}{HTML}{2A78D6}
\definecolor{srOrange}{HTML}{EB6834}
\definecolor{srGreen}{HTML}{1BAF7A}
\definecolor{srYellow}{HTML}{EDA100}
\definecolor{srPurple}{HTML}{6250D6}
\definecolor{srInk}{HTML}{1F1F1D}
\definecolor{srMuted}{HTML}{52514E}
\definecolor{srGrid}{HTML}{D8D7D2}

\newcommand{\srfont}{\sffamily\fontsize{7}{8.2}\selectfont}
\newcommand{\srfontb}{\sffamily\bfseries\fontsize{7}{8.2}\selectfont}

\tikzset{
  sr/base/.style={draw=srInk!75, line width=0.5pt, rounded corners=1.6pt, align=center,
    inner xsep=2.5pt, inner ysep=2pt, font=\srfont, text=srInk, minimum height=9.5mm},
  % client and edge (browser, CloudFront, S3)
  sr/client/.style={sr/base, fill=srBlue!10},
  % identity, scheduling, operations (Cognito, EventBridge, CloudWatch, SES)
  sr/service/.style={sr/base, fill=srPurple!8},
  % compute (Lambda functions, API Gateway)
  sr/compute/.style={sr/base, fill=srOrange!11},
  % queues (SQS)
  sr/queue/.style={sr/base, fill=srYellow!16},
  % data (DynamoDB)
  sr/data/.style={sr/base, fill=srGreen!11},
  % numbered step marker, placed on a path: node[sr/step]{3}
  sr/step/.style={circle, fill=srInk, text=white, font=\sffamily\bfseries\fontsize{6}{6}\selectfont,
    inner sep=0pt, minimum size=8.4pt},
  % arrows
  sr/flow/.style={draw=srInk, line width=0.55pt, -{Stealth[length=3.6pt,width=2.8pt]}},
  sr/async/.style={sr/flow, dash pattern=on 2.4pt off 1.4pt},
  sr/reply/.style={sr/flow, densely dotted},
  sr/aux/.style={draw=srMuted, line width=0.45pt, dotted, -{Stealth[length=3.2pt,width=2.6pt]}},
  % arrow and annotation labels
  sr/lbl/.style={font=\srfont, text=srInk, align=center, inner sep=1pt},
  sr/note/.style={font=\srfont, text=srMuted, align=center, inner sep=1pt},
}
. Size: about 15.9 cm x 5.6 cm. Not drawn to time scale.
%
% Suggested caption:
%   Request sequences (not to scale). (a) A command is accepted (202) as soon as its CMD record
%   is written and it is enqueued; the browser polls GET /commands/{id} (300, 500, 1000, 2000 ms
%   back-off) until the worker's transaction has settled it. Server settle latency =
%   settledAt - acceptedAt, both stored on the CMD record; the client-observed latency adds at
%   most one poll interval. (b) Each offer schedules its own timer on the standard delay queue;
%   at expiresAt the forwarder moves the TICK into the workshop's FIFO group, and the worker
%   expires the offer (tick-first) and relays the seat in the same decision. Backfill latency =
%   reallocatedAt - releasedAt, where releasedAt = expiresAt for an expiry.
\begin{tikzpicture}[x=1cm, y=1cm]
  \def\lifeline#1#2{\draw[srMuted!70, line width=0.4pt, dash pattern=on 1.6pt off 1.2pt] (#1,-0.28) -- (#1,#2);}
  \def\head#1#2#3#4{\node[#3, minimum height=0.48cm, minimum width=1.05cm, inner ysep=1pt] at (#1,#2) {#4};}
  % message labels are drawn on white so that a lifeline they cross does not run through the text
  \def\msg#1#2#3#4{\draw[sr/flow] (#1,#3) -- node[sr/lbl, fill=white, inner sep=0.5pt, above=0.6pt] {#4} (#2,#3);}
  \def\rmsg#1#2#3#4{\draw[sr/reply] (#1,#3) -- node[sr/lbl, fill=white, inner sep=0.5pt, above=0.6pt] {#4} (#2,#3);}
  \def\amsg#1#2#3#4{\draw[sr/async] (#1,#3) -- node[sr/lbl, fill=white, inner sep=0.5pt, above=0.6pt] {#4} (#2,#3);}
  \def\actbar#1#2#3#4{\fill[#4] (#1-0.065,#2) rectangle (#1+0.065,#3);}
  \def\swatch#1#2#3{\fill[#3] (#1,#2-0.07) rectangle (#1+0.22,#2+0.07);}
  \def\note#1#2#3#4{\node[sr/base, fill=white, minimum height=0.34cm, inner ysep=1pt, #3] at (#1,#2) {#4};}

  % ===================================================== (a) command path
  \begin{scope}
    \def\C{0.50}\def\A{2.00}\def\D{3.55}\def\Q{5.05}\def\W{6.55}
    \node[sr/lbl, anchor=west, font=\srfontb] at (-0.05,0.55) {(a) command: 202 $\to$ poll $\to$ settled};
    \foreach \x in {\C,\A,\D,\Q,\W} {\lifeline{\x}{-4.62}}
    \actbar{\C}{-0.62}{-3.95}{srBlue!45}
    \actbar{\D}{-1.00}{-2.70}{srGreen!70}
    \actbar{\W}{-2.20}{-2.70}{srOrange!45}
    \head{\C}{0}{sr/client}{Client}
    \head{\A}{0}{sr/compute}{API}
    \head{\D}{0}{sr/data}{CMD table}
    \head{\Q}{0}{sr/queue}{FIFO}
    \head{\W}{0}{sr/compute}{worker}
    \msg{\C}{\A}{-0.62}{POST}
    \msg{\A}{\D}{-1.00}{put ACCEPTED}
    \node[sr/note, anchor=north west] at (\D+0.08,-1.03) {acceptedAt};
    \draw[sr/flow] (\A,-1.46) -- node[sr/lbl, fill=white, inner sep=0.5pt, below=0.6pt, pos=0.80] {SendMessage} (\Q,-1.46);
    \rmsg{\A}{\C}{-1.82}{202 + href}
    \msg{\C}{\A}{-2.24}{GET (poll)}
    \msg{\Q}{\W}{-2.20}{deliver}
    \rmsg{\A}{\C}{-2.64}{ACCEPTED}
    \msg{\W}{\D}{-2.70}{commit: APPLIED}
    \node[sr/note, anchor=north west] at (\D+0.06,-2.74) {settledAt};
    \msg{\C}{\A}{-3.16}{GET (poll)}
    \msg{\A}{\D}{-3.56}{consistent get}
    \rmsg{\A}{\C}{-3.95}{APPLIED}
    % legend: latency definitions
    \swatch{-0.05}{-5.00}{srBlue!45}
    \node[sr/lbl, anchor=west] at (0.22,-5.00) {client-observed: POST $\to$ terminal poll response};
    \swatch{-0.05}{-5.35}{srGreen!70}
    \node[sr/lbl, anchor=west] at (0.22,-5.35) {server settle $=$ settledAt $-$ acceptedAt (CMD record)};
  \end{scope}

  % ===================================================== (b) expiry relay
  \begin{scope}[xshift=7.75cm]
    \def\WA{0.62}\def\D{2.25}\def\F{3.85}\def\Q{5.30}\def\WB{6.80}
    \node[sr/lbl, anchor=west, font=\srfontb] at (-0.05,0.55) {(b) offer $\to$ timer $\to$ tick $\to$ backfill};
    \foreach \x in {\WA,\D,\F,\Q,\WB} {\lifeline{\x}{-4.62}}
    \actbar{\WB}{-2.52}{-3.62}{srOrange!70}
    \head{\WA}{0}{sr/compute}{worker}
    \head{\D}{0}{sr/queue}{expiry-delay}
    \head{\F}{0}{sr/compute}{forwarder}
    \head{\Q}{0}{sr/queue}{FIFO}
    \head{\WB}{0}{sr/compute}{worker}
    \note{\WA-0.55}{-0.62}{anchor=west}{decide(): offer O to A}
    \amsg{\WA}{\D}{-1.14}{TICK(O)}
    \node[sr/note, anchor=north] at ($(\WA,-1.16)!0.5!(\D,-1.16)$) {delay $= \mathit{hold}$};
    \note{\WA-0.55}{-1.72}{anchor=west}{commit offer O}
    % the hold elapses (break)
    \node[sr/note, fill=white, inner sep=0.4pt] at (3.70,-2.06) {$\cdots$ hold elapses (60\,s in the experiments, default 900\,s) $\cdots$};
    \draw[srMuted, line width=0.4pt, dash pattern=on 3pt off 1.5pt] (-0.05,-2.52) -- (7.30,-2.52);
    \node[sr/note, anchor=south east, inner sep=0.6pt] at (7.30,-2.50) {expiresAt $=$ releasedAt};
    \amsg{\D}{\F}{-2.86}{TICK due}
    \amsg{\F}{\Q}{-3.20}{TICK, group}
    \amsg{\Q}{\WB}{-3.50}{deliver}
    \note{\WB+0.10}{-3.90}{anchor=east, align=right}{tick-first: O expired;\\O2 offered to B at reallocatedAt}
    \amsg{\WB}{\D}{-4.50}{TICK(O2) for the next hold}
    % legend
    \swatch{-0.05}{-5.00}{srOrange!70}
    \node[sr/lbl, anchor=west] at (0.22,-5.00) {backfill latency $=$ reallocatedAt $-$ releasedAt};
    \node[sr/lbl, anchor=west] at (0.22,-5.35) {a late or duplicate TICK finds nothing due (no-op)};
  \end{scope}
\end{tikzpicture}

  \caption{Request sequences (not to scale). (a)~A command is accepted (202) once its record is
  written and it is enqueued; the browser polls until the worker's transaction has settled it.
  Server settle latency = \code{settledAt} $-$ \code{acceptedAt}; the client-observed latency adds
  at most one poll interval. (b)~Each offer schedules its own timer; at \code{expiresAt} the forwarder
  moves the TICK into the workshop's FIFO group, and the worker expires the offer (tick-first) and
  relays the seat in the same decision. Backfill latency = \code{reallocatedAt} $-$
  \code{releasedAt}, with \code{releasedAt} = \code{expiresAt} for an expiry.}
  \label{fig:seq}
\end{figure}

\subsection{Availability, security and privacy}
Within the region every managed service is multi-AZ, so an AZ failure needs no action. A worker
crash delays only its own workshop until the message is redelivered (120\,s visibility timeout);
other workshops continue. A region failure is not covered: recovery means restoring the table from
PITR, which keeps recovery points for up to 35 days at per-second granularity~\cite{aws-ddb-pitr}
with the latest restorable time typically five minutes in the past~\cite{aws-ddb-pitr-howitworks}
(recovery point objective of about five minutes), and redeploying the stacks; we have not measured
that recovery time and claim no recovery time objective. The actor of every command comes only from
the verified token, all input is schema-validated, organiser rights are checked in the API and again
in the worker (C7), and each function has its own least-privilege role (e.g.\ the sweep can only read the table, the
forwarder cannot access it). The API is throttled per stage, cross-origin requests are admitted only
from the SPA's origins, and the S3 bucket is private. We store only e-mail address and display name,
never log raw request events, and keep logs for 14 days. E-mail goes through SES,
which in its sandbox delivers only to verified addresses~\cite{aws-ses-production-access}; until
production access is granted, pilot participants rely on in-app notifications.
